\section{}

\paragraph{1275年}\eventanchor{SS-1275-REGNAL-YEAR}{SS-1275-REGNAL-YEAR}　南宋以宋恭帝在位第1年作为诏令、赋役与外交文书的纪年框架。账本以《宋史》卷24—47〈本纪〉；《建炎以来系年要录》本年条；《宋会要辑稿》相关门类标明来源路径；该锚点连接编年节点，关于数量、实施和地方后果仍应遵守材料的文体与立场限制。

\paragraph{1276年}\eventanchor{SS-1276-EVENT-148}{SS-1276-EVENT-148}　元军进入临安，宋恭帝降元。账本以《宋史》卷24—47〈本纪〉；《建炎以来系年要录》本年条；《宋会要辑稿》相关门类标明来源路径；该锚点连接编年节点，关于数量、实施和地方后果仍应遵守材料的文体与立场限制。

\paragraph{1276年}\eventanchor{SS-1276-REGNAL-YEAR}{SS-1276-REGNAL-YEAR}　南宋以宋恭帝在位第2年作为诏令、赋役与外交文书的纪年框架。账本以《宋史》卷24—47〈本纪〉；《建炎以来系年要录》本年条；《宋会要辑稿》相关门类标明来源路径；该锚点连接编年节点，关于数量、实施和地方后果仍应遵守材料的文体与立场限制。

\paragraph{1277年}\eventanchor{SS-1277-EVENT-149}{SS-1277-EVENT-149}　南宋遗臣、宗室和地方军民在东南继续抵抗。账本以《宋史》卷24—47〈本纪〉；《建炎以来系年要录》本年条；《宋会要辑稿》相关门类标明来源路径；该锚点连接编年节点，关于数量、实施和地方后果仍应遵守材料的文体与立场限制。

\paragraph{1277年}\eventanchor{SS-1277-REGNAL-YEAR}{SS-1277-REGNAL-YEAR}　南宋遗民政权以宋端宗在位第1年作为诏令、赋役与外交文书的纪年框架。账本以《宋史》卷24—47〈本纪〉；《建炎以来系年要录》本年条；《宋会要辑稿》相关门类标明来源路径；该锚点连接编年节点，关于数量、实施和地方后果仍应遵守材料的文体与立场限制。

\paragraph{1278年}\eventanchor{SS-1278-REGNAL-YEAR}{SS-1278-REGNAL-YEAR}　南宋遗民政权以宋端宗在位第2年作为诏令、赋役与外交文书的纪年框架。账本以《宋史》卷24—47〈本纪〉；《建炎以来系年要录》本年条；《宋会要辑稿》相关门类标明来源路径；该锚点连接编年节点，关于数量、实施和地方后果仍应遵守材料的文体与立场限制。

\paragraph{1279年}\eventanchor{SS-1279-EVENT-150}{SS-1279-EVENT-150}　崖山海战后南宋政权终结。账本以《宋史》卷24—47〈本纪〉；《建炎以来系年要录》本年条；《宋会要辑稿》相关门类标明来源路径；该锚点连接编年节点，关于数量、实施和地方后果仍应遵守材料的文体与立场限制。

\paragraph{1279年}\eventanchor{SS-1279-REGNAL-YEAR}{SS-1279-REGNAL-YEAR}　南宋遗民政权以宋少帝赵昺在位第1年作为诏令、赋役与外交文书的纪年框架。账本以《宋史》卷24—47〈本纪〉；《建炎以来系年要录》本年条；《宋会要辑稿》相关门类标明来源路径；该锚点连接编年节点，关于数量、实施和地方后果仍应遵守材料的文体与立场限制。

